\subsection{交换图}

例如，考虑下面的图：

\[
\begin{array}{c} \mathrm{A} \xrightarrow {f} \mathrm{B} \xrightarrow {g} \mathrm{C} \xrightarrow {h} \mathrm{D} \\ a \Big \downarrow \qquad b \Big \downarrow \qquad c \Big \downarrow \qquad d \Big \downarrow \\ \mathrm{A} ^ {\prime} \xrightarrow [ f ^ {\prime} ]{} \mathrm{B} ^ {\prime} \xrightarrow [ g ^ {\prime} ]{} \mathrm{C} ^ {\prime} \xrightarrow [ h ^ {\prime} ]{} \mathrm{D} ^ {\prime} \end{array}\tag{2}
\]

对每条由图中若干线段按箭头所示方向组成的路径，都对应着一个映射，它从第一条线段起点所表示的集合映到最后一条线段终点所表示的集合，即所经各线段所表示的映射的复合。对图的每个顶点，例如 B，按约定有一条化为 B 自身的路径，与之对应的是恒等映射 \(l_{B}\) 。

在 (2) 中，例如有三条从 A 出发到 C' 终止的路径；相应的映射为 \(c \circ g \circ f\) 、 \(g' \circ b \circ f\) 和 \(g' \circ f' \circ a\) 。称一个图为交换的，如果对图中每一对具有相同起点与终点的路径，相应的两个映射都相等；特别地，若一条路径的起点与终点重合，则相应的映射必须是恒等映射。

为使图 (2) 交换，必要且充分的是关系式

\[
f ^ {\prime} \circ a = b \circ f, \quad g ^ {\prime} \circ b = c \circ g, \quad h ^ {\prime} \circ f = d \circ h;\tag{3}
\]

成立；换言之，必要且充分的是 (2) 中所含的三个方形图都是交换的。因为关系式 (3) 蕴含 \(c \circ g \circ f = g' \circ b \circ f\) （由于 \(c \circ g = g' \circ b\) ），以及 \(g' \circ b \circ f = g' \circ f' \circ a\) （由于 \(b \circ f = f' \circ a\) ）；于是从 A 出发到 C' 终止的三条路径给出同一映射。类似地可以验证，从 A 出发到 D' 终止的四条路径（相应地，从 B 出发到 D' 终止的三条路径）给出同一映射。关系式 (3) 表示：从 A（相应地 B, C）出发到 B'（相应地 C', D'）终止的两条路径给出同一映射。(2) 的其他顶点对都不能用多于一条路径连接，因此图 (2) 是交换的。

在下文中，我们把其他类型图形的类似结果留给读者验证。

